\appendix
\section{T3 hint-ablation results}
\label{app:hint-ablation-table}

\begin{table*}[!t]
\centering
\small
\caption{3-column hint study on T3 cross-function inequalities
(pass@$5$). The \textsc{Sub+Chain} chain was reverse-engineered
from an Aristotle T3 closure, so this column is a
proof-plan-condition upper bound, not an independent capability
test. Generalist closure rises from $0$--$1/6$ at \textsc{Base}
to $3$--$5/6$ at \textsc{Sub+Chain}. DSPv2 closes the same
single target across both hint columns. Footnote superscripts
match Table~\ref{tab:headline}.}
\label{tab:ablation}
\begin{tabular}{lccc}
\toprule
Drafter & \textsc{Base} & \textsc{Sub} & \textsc{Sub+Chain} \\
\midrule
Claude Sonnet~4.6  & 0/6 & 1/6 & 4/6 \\
Kimi~K2.5          & 1/6 & 1/6 & 3/6 \\
Gemini~3 Pro       & 0/6 & 2/6 & 4/6 \\
Mistral Large~3$^{\ast}$ & 0/6 & 0/6 & 0/6 \\
\midrule
Claude Opus~4.6$^{\ddagger\ddagger}$ & 0/6 & 2/6 & 5/6 \\
DSPv2-$7$B GPTQ-int$8$ & 0/6 & 1/6 & 1/6 \\
\bottomrule
\end{tabular}
\end{table*}

\section{Hint-ablation legend table}
\label{app:hint-legend-table}

\begin{table*}[!t]
\centering
\small
\setlength{\tabcolsep}{4pt}
\caption{Hint-ablation legend, boxed-table form. Same content as the
inline legend in~\S\ref{sec:methods-hint-ablation}.}
\label{tab:hint-legend}
\begin{tabular}{@{}lp{0.78\textwidth}@{}}
\toprule
\textsc{Base}      & Goal statement only, no hint (no-hint baseline). \\
\textsc{Sub}       & \textsc{Base} plus $\Phi$-substitution hint
                     $x \mapsto \exp(\lambda^{\star} x -
                     (\lambda^{\star})^{2}\sigma^{2}/2)$,
                     family-specific $\lambda^{\star}$ (lemma search
                     given the substitution). \\
\textsc{Sub+Chain} & \textsc{Sub} plus an explicit Mathlib lemma chain
                     (\texttt{Real.exp\_le\_one\_iff},
                     \texttt{Real.log\_nonneg},
                     \texttt{Real.sqrt\_le\_sqrt}, \dots)
                     (lemma application given both). \\
\bottomrule
\end{tabular}

\end{table*}

\section{Benchmark format}
\label{app:bench-format-fig}

The two listings below show one T3 target as it appears in the Lean
source (Listing~\ref{lst:bench-lean}) and as a row of
\texttt{benchmark.jsonl} (Listing~\ref{lst:bench-json}). The full
field-level schema is in Appendix~\ref{app:schema}.

\begin{lstlisting}[style=leanstyle, caption={Lean theorem statement
  for the T3 target \texttt{etavector\_eq\_sqrt\_two\_mul\_etahr}.
  Stored verbatim in the \texttt{theorem\_statement} field of the
  matching JSONL row.},
  label={lst:bench-lean}]
theorem etavector_eq_sqrt_two_mul_etahr
    (b : Real) (hb : 1 <= b) :
    etaVector b = Real.sqrt 2 * etaHR b := by
  unfold etaVector etaHR
  rw [Real.sqrt_mul]
  . ring
  . positivity
\end{lstlisting}

\begin{lstlisting}[style=jsonstyle, caption={\texttt{benchmark.jsonl}
  record (excerpt) for the same target. The
  \texttt{hint\_ablation} block carries the three prompt columns
  for T3 targets only (Table~\ref{tab:hint-legend}).},
  label={lst:bench-json}]
{
  "target_id": "etavector_eq_sqrt_two_mul_etahr",
  "tier": "T3",
  "family": "cross",
  "requires_formal_avs_lean": true,
  "required_imports": [
    "FormalAvs.HowardRamdas",
    "FormalAvs.WhitehouseVector"
  ],
  "closure_rate_n5": {
    "sonnet_4_6":  0.0,
    "kimi_k2_5":   0.6,
    "gemini_3pro": 0.0,
    "mistral_l3":  0.4
  },
  "hint_ablation": {
    "BASE":      {"sonnet_4_6": 0.0, ...},
    "SUB":       {"sonnet_4_6": 0.2, ...},
    "SUB+CHAIN": {"sonnet_4_6": 0.8, ...}
  },
  "aristotle_session_id": "74fe5c49"
}
\end{lstlisting}

\section{Evaluation-protocol pipeline}
\label{app:eval-protocol-fig}

Table~\ref{tab:eval-protocol} expands the five-stage drafter sweep
pipeline summarised in~\S\ref{sec:methods-hint-ablation}. Every
(target, drafter, prompt) cell traverses the five stages in order.
Stage~4 pins outcomes to the fixed logical foundation
of~\S\ref{sec:methods-qa}.

\begin{table*}[!t]
\centering
\small
\setlength{\tabcolsep}{4pt}
\caption{Five-stage evaluation pipeline applied per (target,
drafter, prompt) cell. Companion to
\S\ref{sec:methods-hint-ablation}. Expanded form of the protocol
text in~\S\ref{sec:methods-tiers}.}
\label{tab:eval-protocol}
\begin{tabular}{@{}rlp{0.65\textwidth}@{}}
\toprule
\# & Stage & Action \\
\midrule
1 & Target  & Lean theorem read from \texttt{benchmark.jsonl}
              ($60$ targets), tagged with tier T0--T5 and required
              imports. \\
2 & Prompt  & Hint-list or neutral template. For T3, one of
              \textsc{Base} / \textsc{Sub} / \textsc{Sub+Chain}
              (Table~\ref{tab:hint-legend}). \\
3 & Sample  & $N{=}5$ attempts per cell: $1$ greedy at $T{=}0$
              plus $4$ samples at $T{=}0.7$, $8$k tokens each. \\
4 & Verify  & \texttt{lake build} on the candidate, reject any
              \texttt{sorry}, then assert
              $\texttt{\#print axioms} \subseteq
              \mathcal{A}_{\text{core}}$. \\
5 & Score   & Target closed for the cell iff $\geq 1$ of $N$
              attempts passes stage~4. Aggregate as
              per-cell pass@$N$. \\
\bottomrule
\end{tabular}

\end{table*}

\section{Dataset schema}
\label{app:schema}

Each row of \texttt{benchmark.jsonl} contains the following
fields.

\begin{itemize}[leftmargin=1.3em, itemsep=0.15em]
\item \texttt{target\_id}: unique identifier (stable across versions).
\item \texttt{theorem\_statement}: the Lean~4 theorem statement as a
raw string (imports included).
\item \texttt{tier}: one of \texttt{T0}, \texttt{T1}, \texttt{T2},
\texttt{T3}, \texttt{T4}, \texttt{T5}.
\item \texttt{requires\_formal\_avs\_lean}: boolean flag. True if
the statement references a \texttt{formal-avs-lean} definition
not in Mathlib.
\item \texttt{family}: one of \texttt{hr}, \texttt{vector},
\texttt{betting}, \texttt{aCS}, \texttt{cross} (cross-family
inequality).
\item \texttt{closure\_rate\_n5}: map from drafter-name to fraction
closed at $N=5$.
\item \texttt{hint\_ablation}: present for T3 targets. Map from
hint-column key (\texttt{"BASE"}, \texttt{"SUB"}, or
\texttt{"SUB+CHAIN"} as in Table~\ref{tab:hint-legend}) to
closure-rate map across drafters.
\item \texttt{aristotle\_session\_id}: fingerprint hash linking to a
row in \texttt{aristotle\_history.jsonl} (if present).
\item \texttt{mathlib\_lemmas}: for T2/T3, the Mathlib lemma chain
that admits the closure (used in the \textsc{Sub+Chain} prompt).
\end{itemize}

\section{Aristotle history schema}
\label{app:aristotle}

Each row of \texttt{aristotle\_history.jsonl} contains the following
fields.

\begin{itemize}[leftmargin=1.3em, itemsep=0.15em]
\item \texttt{session\_id}: short fingerprint (8-12 hex chars).
\item \texttt{target\_id}: pointer into \texttt{benchmark.jsonl}.
\item \texttt{submitted\_at}, \texttt{completed\_at}: ISO-8601 UTC
timestamps.
\item \texttt{verdict}: one of \texttt{closed}, \texttt{partial},
\texttt{refuted}.
\item \texttt{axiom\_audit}: sorted list of axioms used by the closed
Lean kernel, as reported by \texttt{\#print axioms}.
\item \texttt{counterexample}: present only when \texttt{verdict =
refuted}. The Lean-syntax counterexample witness.
\item \texttt{residual\_sorries}: list of \texttt{sorry} names present
in the partial proof (present only when \texttt{verdict = partial}).
\item \texttt{compute\_minutes}: elapsed wall-clock time of the Aristotle session.
\end{itemize}

\section{Drafter prompts}
\label{app:prompts}

\subsection{Generalist drafter prompt (closed-source)}

\begin{lstlisting}[basicstyle=\ttfamily\small,breaklines=true,breakautoindent=false,columns=flexible]
You are a Lean 4 expert proficient with Mathlib.
Close the following theorem with a clean tactic proof.
Prefer concise tactic chains. Avoid `sorry`. Avoid axioms.

<theorem-statement>

<optional hint block>
\end{lstlisting}

The \textsc{Sub} hint adds: ``Hint: the log-domain substitution
$x \mapsto \exp(\lambda^\star x - (\lambda^\star)^2 \sigma^2 / 2)$ with
family-specific $\lambda^\star$ is likely useful.''

The \textsc{Sub+Chain} hint adds on top: ``Hint: the following Mathlib lemmas
are likely relevant, in approximate order: \texttt{Real.exp\_le\_one\_iff},
\texttt{Real.log\_nonneg}, \texttt{Real.sqrt\_le\_sqrt},
\texttt{Real.exp\_log\_eq\_iff}, \texttt{Real.sqrt\_mul\_self},
\texttt{field\_simp}, \texttt{positivity}.''

\subsection{Specialised prover prompts (DSPv2, Goedel-V2)}

DSPv2 is autocomplete-trained: the input terminates at the trailing
\texttt{:=} token and the model emits the proof body directly.

\begin{lstlisting}[basicstyle=\ttfamily\small,breaklines=true,breakautoindent=false,columns=flexible]
<imports>

<opens>

namespace <scratch_namespace>

<theorem-header> :=
\end{lstlisting}

Goedel-V2 is chat-trained: per the model card, the input is wrapped
as a single user turn that asks for a proof plan followed by the
final code in a fenced \texttt{lean4} block. We extract the proof
body from that block and verify against the same lake project as
the rest of the sweep.

\begin{lstlisting}[basicstyle=\ttfamily\small,breaklines=true,breakautoindent=false,columns=flexible]
Complete the following Lean 4 code:

```lean4
<imports>

namespace <scratch_namespace>

<theorem-header> := by
  sorry
```

Before producing the Lean 4 code to formally prove the given
theorem, provide a detailed proof plan ...
\end{lstlisting}

The \textsc{Sub+Chain} column adds the Mathlib-lemma hint as a Lean
\texttt{/- hint -/} block before the trailing \texttt{:=} token (DSPv2)
or in the user turn before the code block (Goedel-V2). DSPv2 is
prefix-only and ignores non-prefix hint context. We retain all three
hint columns for both rows for parity with the closed-source
generalists.

We use community quantizations: DSPv2 GPTQ-int8 (one of the 15
community variants linked from the official model page), Goedel-V2
GGUF Q$6$\_K. Both are inference-only artifacts. The underlying
weights are the official 7B / 8B releases.

\subsection{Alternative T3 hint chain from library reference proofs}
\label{app:alternative-hint-chain}

To rule out the concern that the \textsc{Sub+Chain} hint contaminates
the ablation by replaying Aristotle's tactic-prior, we extracted
a second hint chain directly from the closed reference proofs in
\texttt{formal-avs-lean} (the per-target proofs the library
authors hand-wrote for the T3 statements, prior to any solver
run). The two chains are reported below:

\begin{sloppypar}
\begin{description}[leftmargin=1.3em, itemsep=0.2em]
\item[\textsc{Sub+Chain} headline.]
\texttt{Real.exp\_le\_one\_iff}, \texttt{Real.log\_nonneg},
\texttt{Real.sqrt\_le\_sqrt}, \texttt{Real.exp\_log\_eq\_iff},
\texttt{Real.sqrt\_mul\_self}, \texttt{field\_simp},
\texttt{positivity}.
\item[Library reference proofs.]
\texttt{Real.sqrt\_le\_sqrt}, \texttt{Real.sqrt\_mul},
\texttt{Real.log\_nonneg}, \texttt{Real.le\_sqrt\_of\_sq\_le},
\texttt{nlinarith}, \texttt{field\_simp},
\texttt{Real.log\_two\_gt\_d9}.
\end{description}
\end{sloppypar}

The two chains overlap on $5/7$ lemmas
(\texttt{Real.sqrt\_le\_sqrt}, \texttt{Real.log\_nonneg},
\texttt{field\_simp}, \texttt{positivity} or
\texttt{Real.le\_sqrt\_of\_sq\_le}, and a sqrt-multiplication
identity). The non-overlapping
parts reflect different proof shapes (the \textsc{Sub+Chain} chain
includes \texttt{Real.exp\_log\_eq\_iff} for the substitution
step. The reference-proof chain uses \texttt{nlinarith} +
\texttt{Real.log\_two\_gt\_d9} for numerical bounds). Both
chains live in the same Mathlib lemma neighbourhood. The
non-overlap reflects different proof shapes, not distribution
shift.

The headline ablation in Table~\ref{tab:ablation} uses the
\textsc{Sub+Chain} chain. The library-reference chain is released in
the supplementary archive as an alternative hint set for
researchers wanting an Aristotle-independent baseline.

\paragraph{Spot-check N=5 result.}
We ran a small Sonnet $N=5$ comparison on $3$ of the $6$ T3
targets (\texttt{etabetting\_le\_etahr},
\texttt{etavector\_eq\_sqrt\_two\_mul\_etahr},
\texttt{etahr\_le\_etavector}) using the library-reference
chain. Sonnet closed $0/15$ attempts, versus the headline \textsc{Sub+Chain}
result of $4/6$ across all six T3 targets. The result is
informative: the specific lemma sequence in the \textsc{Sub+Chain} chain
is consequential within the broad ``Mathlib lemma chain'' hint
category. The two chains overlap on $5/7$ lemmas but produce
different per-target close-rates at small $N$, suggesting
that hint composition (lemma order, presence of specific bridging
identities like \texttt{Real.exp\_log\_eq\_iff}) is itself a
capability axis. A full $4 \times 6 \times N=5$ comparison
across all four generalist drafters and all six T3 targets is
deferred to a future revision.

\section{Tier-distribution table}
\label{app:tier-dist}

\begin{table*}[!t]
\centering
\small
\caption{Target distribution across tier $\times$ family. $60$
targets total. T4 plus T5 are the empirical failure categories
($5$ theorems). $53$ of $60$ non-frontier targets require
\texttt{formal-avs-lean}.}
\label{tab:tier-dist}
\begin{tabular}{lcccccc}
\toprule
Family & T0 & T1 & T2 & T3 & T4 & T5 \\
\midrule
HR-family                    & 2 & 11 & 3 & - & - & - \\
betting family               & 1 & 10 & 1 & - & - & - \\
Whitehouse vector            & 1 &  6 & 2 & - & - & - \\
Asymptotic CLT               & 1 &  7 & 2 & - & - & - \\
Cross-family (inequality)    & - & - & - & 6 & - & - \\
Mathlib-only reference       & 2 & - & - & - & - & - \\
Library-gap (Mathlib PR)     & - & - & - & - & 4 & - \\
Capability-ceiling frontier  & - & - & - & - & - & 1 \\
\midrule
Total                        & 7 & 34 & 8 & 6 & 4 & 1 \\
\bottomrule
\end{tabular}

\end{table*}

\section{Baseline reproducibility}
\label{app:reproducibility}

To reproduce the headline sweep:
\begin{enumerate}[leftmargin=1.3em, itemsep=0.15em]
\item Clone the companion Lean~4 library at the pinned release
commit (snapshot URL in supplementary
\texttt{REPRODUCE.md}, item~\textsc{lib-snapshot}. Lean
toolchain and Mathlib commit pin in \texttt{lean-toolchain}
and \texttt{lakefile.lean}).
\item Clone the dataset repo from the snapshot Hugging Face
location (URL in supplementary \texttt{REPRODUCE.md},
item~\textsc{dataset-snapshot}).
\item Run the provided \texttt{sweep.py} against a drafter
endpoint. It submits each of the $48$ evaluated targets under
the mixed-temperature pass@$5$ schedule (one greedy sample at
$T=0$ followed by four independent samples at $T=0.7$) and
writes per-target closure records. The driver prints the
SHA-256 of the prompt files and dataset snapshot it used.
\item Verify each closure by running \texttt{lake build}
against the produced Lean file. The per-target axiom audit is
extracted via \texttt{lean -{}-run axiom\_audit.lean}.
\item Run \texttt{bash reproduce.sh} from the supplementary
archive to regenerate every numeric cell in
Table~\ref{tab:headline}, Table~\ref{tab:ablation},
Table~\ref{tab:prompt-sensitivity}, and the Wilson-CI tables
from cached drafter outputs (released alongside the dataset).
\end{enumerate}
For the specialised-prover column, the pinned GPTQ-quantised
DSPv2 weights serve as the baseline endpoint.

\paragraph{Asset licenses.}
Benchmark JSONL, tier metadata, per-drafter closure records at
$N=5$, drafter prompts, 3-column hint ablation records, and the
14-session Aristotle history archive are released under CC-BY-4.0.
The underlying Lean~4 library is released under Apache-2.0. The
GPTQ-quantised DSPv2 weights are released as an artifact. Anonymised
asset URLs are included in the supplementary archive.

\paragraph{Wilson 95\% confidence intervals on Table~\ref{tab:headline}.}
Table~\ref{tab:wilson-ci} reports the Wilson 95\% confidence interval
on every per-cell pass@5 closure rate from the headline matrix. The
effective $N$ for each Wilson interval is the per-tier denominator
because pass@5 collapses each target's $5$ attempts into a single
Bernoulli outcome (target-closed if $\geq 1$ of $5$ attempts closed).
The intervals are widest on the $n=1$ T5 cell and the small-$n$ T0,
T2, T3, T4 cells (all per-tier denominator $\leq 8$). The headline
findings are stable under the Wilson CIs. The four-drafter T3 wall
has every drafter $\leq 17\%$ closure CI upper bound, separated from
Aristotle's $6/6$ closure on the same column. Mistral's T1 collapse
$1/34$ has CI $[1, 15]\%$, below the $34$ to $66\%$ band of the
other three drafters. The four-drafter T0 floor at $0/7$ has Wilson
upper bound $35\%$ across all four, consistent with the Mathlib-canon
framing.

\begin{table*}[!t]
\caption{Wilson 95\% confidence intervals on every per-cell
pass@5 closure rate from Table~\ref{tab:headline}. Cells report
targets-closed/tier-size with the Wilson 95\% CI in brackets
(percent units). Mistral T1 $1/34$ at $[1,15]$ is separated
from the $34$--$66\%$ band of Sonnet/Kimi/Gemini. The
four-drafter T3 wall has every drafter at most $17\%$ upper
bound, below Aristotle's $6/6$ T3 closure
(\S\ref{sec:capability-ladder}). T4 and T5 cells report $0$-imputed
outcomes. Intervals reflect sampling uncertainty on the
structural zero.}
\label{tab:wilson-ci}
\centering
\small
\resizebox{\linewidth}{!}{%
\begin{tabular}{l|cccccc}
\toprule
Drafter & T0 & T1 & T2 & T3 & T4 & T5 \\
& ($n=7$) & ($n=34$) & ($n=8$) & ($n=6$) & ($n=4$) & ($n=1$) \\
\midrule
Claude Sonnet~4.6 & 0/7 \textcolor{gray}{[0,35]} & 18/34 \textcolor{gray}{[37,69]} & 0/8 \textcolor{gray}{[0,32]} & 0/6 \textcolor{gray}{[0,39]} & 0/4 \textcolor{gray}{[0,49]} & 0/1 \textcolor{gray}{[0,79]} \\
Kimi~K2.5 & 0/7 \textcolor{gray}{[0,35]} & 17/34 \textcolor{gray}{[34,66]} & 0/8 \textcolor{gray}{[0,32]} & 1/6 \textcolor{gray}{[3,56]} & 0/4 \textcolor{gray}{[0,49]} & 0/1 \textcolor{gray}{[0,79]} \\
Gemini~3 Pro & 0/7 \textcolor{gray}{[0,35]} & 22/34 \textcolor{gray}{[48,79]} & 0/8 \textcolor{gray}{[0,32]} & 0/6 \textcolor{gray}{[0,39]} & 0/4 \textcolor{gray}{[0,49]} & 0/1 \textcolor{gray}{[0,79]} \\
Mistral Large~3 & 0/7 \textcolor{gray}{[0,35]} & 1/34 \textcolor{gray}{[1,15]} & 0/8 \textcolor{gray}{[0,32]} & 1/6 \textcolor{gray}{[3,56]} & 0/4 \textcolor{gray}{[0,49]} & 0/1 \textcolor{gray}{[0,79]} \\
Claude Opus~4.6 & 0/7 \textcolor{gray}{[0,35]} & 13/34 \textcolor{gray}{[24,55]} & 0/8 \textcolor{gray}{[0,32]} & 0/6 \textcolor{gray}{[0,39]} & 0/4 \textcolor{gray}{[0,49]} & 0/1 \textcolor{gray}{[0,79]} \\
DSPv2-$7$B GPTQ-int$8$ & 3/7 \textcolor{gray}{[16,75]} & 8/34 \textcolor{gray}{[12,40]} & 0/8 \textcolor{gray}{[0,32]} & 1/6 \textcolor{gray}{[3,56]} & 0/4 \textcolor{gray}{[0,49]} & 0/1 \textcolor{gray}{[0,79]} \\
\bottomrule
\end{tabular}}
\end{table*}

\subsection{$N=10$ extension}
\label{app:robustness-n10}

To stress the headline closure pattern under a higher per-cell
attempt budget, we extended all six drafter sweeps from $N=5$ to
$N=10$ by firing five additional independent $T=0.7$ samples per
(target, drafter) cell, for a total of one greedy plus nine $T=0.7$
attempts per cell. The slate, prompts, verification regime, and
v2-strict sorry-detection are held identical to the $N=5$ headline.
The four closed-source drafters (Sonnet, Kimi, Gemini, Mistral),
the supplementary fifth row (Opus), and the specialised prover
(DSPv2-$7$B GPTQ-int8) were each extended on the full $48$-target
slate.

Across the $6 \times 48 = 288$ (drafter, target) cells, doubling
the attempt budget shifts only \emph{three} target-level closures:
Gemini closes one new T2 target
(\texttt{dichotomy\_\allowbreak universal\_\allowbreak monotonicity\_\allowbreak impossible}), Opus
closes one new T3 target (\texttt{etahr\_over\_etabetting\_gt\_one}),
and DSPv2 closes one new T3 target (\texttt{etahr\_le\_etavector}).
All three new closures are axiom-clean and pass v2-strict (no
\texttt{sorry} in the candidate body). The remaining $285$ cells
reproduce their $N=5$ verdict.

Table~\ref{tab:wilson-ci-n10} reports the per-cell Wilson 95\% CIs
at $N=10$. Intervals widen only for cells whose target-level outcome
changed. The $285$ unchanged cells retain identical intervals to
their $N=5$ counterparts in Table~\ref{tab:wilson-ci}. The headline
findings hold under the $N=10$ stress. Mistral T1 collapse remains
at $1/34$ with CI $[1, 15]\%$. The four-drafter T3 wall has every
drafter at most $17\%$ upper bound (Kimi, Mistral, and Opus all at
$1/6 = 17\%$ point estimate, with Sonnet and Gemini at $0/6$). The
four-drafter T0 floor remains at $0/7$ for all four with Wilson
upper bound $35\%$.

\begin{table*}[!t]
\centering
\small
\caption{Per-cell pass@10 closure rates with Wilson 95\% CIs.
T4 and T5 omitted (no drafter closes either at any $N$). Three
cells differ from pass@5 in Table~\ref{tab:wilson-ci}: Gemini T2
$0/8 \to 1/8$, Opus T3 $0/6 \to 1/6$, DSPv2 T3 $1/6 \to 2/6$.
Per-cell intervals are stable across $N=5 \to N=10$.}
\label{tab:wilson-ci-n10}
\begin{tabular}{l|cccc}
\toprule
Drafter & T0 & T1 & T2 & T3 \\
& ($n=7$) & ($n=34$) & ($n=8$) & ($n=6$) \\
\midrule
Claude Sonnet~4.6 & 0/7 \textcolor{gray}{[0,35]} & 18/34 \textcolor{gray}{[37,69]} & 0/8 \textcolor{gray}{[0,32]} & 0/6 \textcolor{gray}{[0,39]} \\
Kimi~K2.5 & 0/7 \textcolor{gray}{[0,35]} & 17/34 \textcolor{gray}{[34,66]} & 0/8 \textcolor{gray}{[0,32]} & 1/6 \textcolor{gray}{[3,56]} \\
Gemini~3 Pro & 0/7 \textcolor{gray}{[0,35]} & 22/34 \textcolor{gray}{[48,79]} & 1/8 \textcolor{gray}{[2,47]} & 0/6 \textcolor{gray}{[0,39]} \\
Mistral Large~3 & 0/7 \textcolor{gray}{[0,35]} & 1/34 \textcolor{gray}{[1,15]} & 0/8 \textcolor{gray}{[0,32]} & 1/6 \textcolor{gray}{[3,56]} \\
Claude Opus~4.6 & 0/7 \textcolor{gray}{[0,35]} & 13/34 \textcolor{gray}{[24,55]} & 0/8 \textcolor{gray}{[0,32]} & 1/6 \textcolor{gray}{[3,56]} \\
DSPv2-$7$B GPTQ-int$8$ & 3/7 \textcolor{gray}{[16,75]} & 8/34 \textcolor{gray}{[12,40]} & 0/8 \textcolor{gray}{[0,32]} & 2/6 \textcolor{gray}{[10,70]} \\
\bottomrule
\end{tabular}

\end{table*}

\paragraph{DSPv2 inference-path note.}
The DSPv2 $N=10$ extension required switching inference paths
(the original loader was unavailable at extension time). The first
$N=10$ attempt on the substitute path returned candidates with
interior whitespace stripped
(``\texttt{:=byintros\dots rw[quantizeReals]}'' in place of
``\texttt{:= by intros\dots rw [quantizeReals]}''). We diagnosed
this as a chat-style-instruction prompt fed to a Lean-autocompletion
model (DSPv2 is a completion model, not a chat model) and switched
to a Lean-autocomplete prompt format
(\texttt{\{imports\}\textbackslash n\textbackslash n\{theorem\_header\} := }).
Re-running the sweep with the autocomplete prompt produced clean
Lean candidates and one new T3 closure
(\texttt{etahr\_le\_etavector}). A subsequent SDK release adds a
model-family registry that auto-routes recognized Lean prover
identifiers to the autocomplete format and warns loudly when a
chat-style template is supplied for a known completion model.
Researchers reproducing the DSPv2 sweep through the released SDK
will get the autocomplete format by default without any
caller-side ceremony.

\paragraph{Reproducibility note: trace-on-default for sweep launches.}
The drafter sweep drivers stream every \texttt{litellm.completion}
event to a relational event log by default when the standard
event-sink environment variables are set, allowing post-hoc
per-target audit, cost rollup, and run-level dashboards. The
default-on behaviour is implemented via a small wrapper that each
driver's \texttt{main()} entry invokes. Passing \texttt{--no-trace}
on the CLI skips the trace session entirely (useful for offline
runs without a database sink). Researchers reproducing the headline
sweep should rely on this default. Per-attempt JSON logs under
\texttt{logs/} are written unconditionally and serve as the
reconstruction path if a run was launched without the database
sink available.

\section{Refutation case study: Aristotle counterexample}
\label{app:refutation-case}

Session \texttt{d2755ea2} submitted a candidate theorem asserting that
the Gaussian small-ball lower bound holds without a symmetry
hypothesis on the underlying filtration. Aristotle refuted the claim
by constructing a two-step filtration witness in which the asserted
bound fails by a constant factor. The witness was reduced to a
Lean term and verified by \texttt{lake build}. The refuted statement
was corrected by re-introducing the dropped symmetry hypothesis and
re-submitted as session \texttt{54614669}, which closed cleanly. Both
sessions and the counterexample witness are preserved in the
history archive.

\section{Extended capability analysis}
\label{app:extended}

This appendix supplies the four subsections trimmed from the main
results in order to keep the main text within the NeurIPS~2026 page
budget. The findings are supporting evidence for the main-text
headline (Aristotle $6/6$ T3 against drafter aggregate $2/24$). The
main text's §\ref{sec:results} retains the headline findings.

\subsection{Aristotle history archive}
\label{app:aristotle-history}

The companion archive of fourteen Aristotle sessions records:
\begin{itemize}[leftmargin=1.3em, itemsep=0.15em]
\item \textbf{Ten fully-closed sessions} with clean axiom audit
$\{$\texttt{propext}, \texttt{Classical.choice}, \texttt{Quot.sound}$\}$,
covering HR admissibility, betting CS admissibility, wealth-process
martingale property, power-loss bound, Gaussian small-ball lower
bound, the universal $\Phi$-transform admissibility for all four
families (HR, vector, aCS, and implicitly betting via the Kelly-optimality
characterisation), matching sharp constants $c_{\mathrm{HR}}$ and
$c_{\mathrm{vector}}$ with their ranking proof, and two
information-theoretic bounds (quantized entropy and Kelly-optimal
betting).
\item \textbf{Three partial-closure sessions} where Aristotle
closed the headline theorem modulo residual specification sorries
(these are recorded with the explicit residual list).
\item \textbf{One refutation session} where Aristotle returned a
counterexample: the submitted claim required a symmetry hypothesis
that had been dropped during IP-strip. Aristotle's counterexample
witness was preserved as an auxiliary dataset entry and triggered
a correction to the original statement.
\end{itemize}

The 14-session archive supplies ground-truth closures for a subset
of benchmark targets, which allows drafter-side closures to be
cross-validated against Aristotle's certified output.

\subsection{Aristotle capability demonstrations beyond the main table}
\label{app:aristotle-companion}

In addition to the benchmark-Table~\ref{tab:headline} evaluation,
Harmonic Aristotle was exercised on three companion-archive theorems
during library development:
\texttt{c\_sharp\_ranking} (the sharp-constant ranking
$c_{\mathrm{aCS}} \leq c_{\mathrm{vector}}$, a sharp-constant
statement distinct from the rate-function ranking
$\eta_{\mathrm{betting}} \leq \eta_{\mathrm{aCS}} \leq
\eta_{\mathrm{HR}} \leq \eta_{\mathrm{vector}}$ that appears at
T3 of the benchmark),
\texttt{quantized\_entropy\_bound} (an information-theoretic
entropy bound applied to the quantization-coarsened CS), and
\texttt{betting\_is\_kelly\_optimal} (the Kelly-optimality
characterisation of the betting CS). All three closed axiom-clean.
These are \emph{not} Table~\ref{tab:headline} benchmark entries
because the library proofs predate the evaluation. we report them
to contextualise Aristotle's capability on theorems adjacent to the
benchmark targets. No claim about Aristotle's benchmark tier close
rate is made from these companion-archive closures.

\subsection{Cross-drafter tactic-prior diversity}
\label{app:cross-drafter}

On four targets in the benchmark, \emph{every} evaluated closed-source
drafter and Aristotle both close. These dual-closure rows expose
cross-drafter variance: different commercial drafters
close the same theorem via structurally \emph{different tactic chains},
and one drafter may close where another fails at identical attempt
budget. A concrete example is the sharp-constant ranking
\texttt{c\_sharp\_ranking} ($c_{\mathrm{aCS}} \leq c_{\mathrm{vector}}$):
Kimi~K2.5 fails $0/9$ at the \textsc{Sub} hint (its tactic prior directs
search toward \texttt{one\_div\_le\_one\_div\_iff}-style
linear-algebra routes that terminate without discharging the goal),
whereas Gemini~3 Pro closes at att~1 greedy with the concise chain
\texttt{[unfold, gcongr, linarith [Real.pi\_pos]]}. The same target
is also closed by Aristotle under a longer search, and each
of the three closures uses a distinct tactic route. We interpret this
as drafter-specific tactic priors: different drafters have
structurally different views of the Mathlib lemma graph, and the
benchmark's hint-ablation protocol exposes that variance as
information rather than noise.

\paragraph{Convergent T3 close on the $\sqrt{2}$ identity.}
The single T3 target both Kimi and Mistral uniquely close at
$N=5$ greedy is \texttt{etavector\_eq\_sqrt\_two\_mul\_etahr}
($\eta_{\mathrm{vector}} = \sqrt{2}\cdot\eta_{\mathrm{HR}}$). Kimi
closes it on the greedy attempt and reproduces at two $T=0.7$
attempts. Mistral (under the neutral prompt) closes at the greedy
attempt and reproduces at one $T=0.7$ attempt. Sonnet and Gemini
miss this target at $N=5$. Two independent commercial drafters
converging on the same cross-family identity under $0$ explicit
hint is direct evidence that the $\sqrt{2}$ factor between the
HR and Whitehouse-vector rate functions is the most accessible
cross-family relation in the 4-family rate landscape, independent
of drafter tactic-prior.

\subsection{Four solver regimes in the capability-wall tail}
\label{app:four-regimes}

Running the full drafter sweep (Kimi~K2.5, Sonnet~4.6, Gemini~3~Pro)
against the five capability-wall tail theorems under the \textsc{Sub}
and \textsc{Sub+Chain} hint regimes yields $5$ closures in $100$ total
attempts ($5.0\%$ drafter close-rate). Breaking the
closure matrix out by target reveals \emph{four qualitatively
distinct solver regimes} plus a deepest-tier target unreachable by
any generalist drafter (Table~\ref{tab:four-regimes}):

\begin{table*}[!t]
\centering
\caption{Four solver regimes in the no-closure tail. Cells are
(close count)/(attempt count) at $N \leq 9$ per drafter.
$^\dagger$Solver-unreached: no closure across (a) 3 drafters
$\times N=5$ greedy unhinted, (b) 3 drafters $\times$ asymmetric
\textsc{Sub}, (c) 3 drafters $\times N=5$ \textsc{Sub+Chain}, (d) 3
Aristotle attempts. $^\ddagger$Two Aristotle sessions timed out
at $20$~min. One returned a counterexample on the as-stated
form (\S\ref{sec:aristotle-refute}). We avoid claiming
\emph{fundamentally} unreachable.}
\label{tab:four-regimes}
\resizebox{\linewidth}{!}{%
\begin{tabular}{lccccc}
\toprule
Target & Kimi & Sonnet & Gemini & Aristotle & Regime \\
\midrule
\texttt{betting\_is\_kelly\_optimal}      & $1/1$ & $1/1$ & $1/1$ & $1/1$ & \textbf{universal} \\
\texttt{c\_vector\_eq\_sqrt\_two\_mul\_c\_HR} & $1/1$ & $0/5$ & $1/2$ & $1/1$ & \textbf{majority} \\
\texttt{c\_sharp\_ranking}                & $0/9$ & $0/5$ & $1/3$ & $1/1$ & \textbf{drafter-monopoly} \\
\texttt{quantized\_entropy\_bound}        & $0/9$ & $0/5$ & $0/3$ & $1/1$ & \textbf{Aristotle-monopoly} \\
\texttt{equivalence\_\allowbreak break\_at\_\allowbreak finite\_precision} & $0/19$ & $0/15$ & $0/13$ & $0/3$ (timeout$^\ddagger$) & \textbf{solver-unreached$^\dagger$} \\
\bottomrule
\end{tabular}%
}

\end{table*}

The four regimes correspond directly to the capability axes identified
in \S\ref{sec:results}. \emph{Universal} targets exercise only baseline
tactic chains (well within every drafter's Mathlib prior).
\emph{Majority} targets exercise a a rewrite step that most
but not all drafters' priors reach. \emph{Drafter-monopoly} targets
expose cross-drafter tactic-prior divergence (Gemini's \texttt{gcongr}
prior vs Kimi's \texttt{one\_div\_le\_one\_div\_iff} prior).
\emph{Aristotle-monopoly} targets require a multi-step lemma chain
that all generalist drafter priors miss. The
\emph{specialised-only tier} shows the tactic-vocabulary ceiling.
Drafters have the mathematics but lack the Lean~4 API signatures
needed to close. The dataset records the closure matrix at row
granularity, enabling per-target regime analyses.

\subsection{Axiom audit}
\label{app:axiom-audit}

Every closed Lean~4 proof in the release terminates with a
\texttt{\#print axioms} check. The canonical axiom set for a
classical Mathlib proof is
$\mathcal{A}_{\text{core}} \coloneqq \{\texttt{propext},
\texttt{Classical.choice}, \texttt{Quot.sound}\}$. The axiom-set
union across all closed proofs in the release is exactly
$\mathcal{A}_{\text{core}}$, no proof widens the axiom
set beyond this baseline
(Table~\ref{tab:axiom-audit}).

\begin{table*}[!t]
\centering
\small
\caption{Axiom audit. Every closed proof uses exactly
$\mathcal{A}_{\text{core}} = \{\texttt{propext},
\texttt{Classical.choice}, \texttt{Quot.sound}\}$. Refuted
targets return a counterexample witness, no proof. Three
refutations are documented in \S\ref{sec:aristotle-refute} and
\S\ref{app:refutation-case}. No closed proof widens the axiom
set beyond $\mathcal{A}_{\text{core}}$ on either side.}
\label{tab:axiom-audit}
\resizebox{\linewidth}{!}{%
\begin{tabular}{lccc}
\toprule
Source & Verdict & N & Axiom set \\
\midrule
Aristotle main table (T0/T1/T2) & closed & 47 & $\mathcal{A}_{\text{core}}$ \\
Aristotle main table (T1/T2) & refuted & 2 & (counterexample, no proof) \\
Aristotle main table (T3) & closed & 6 & $\mathcal{A}_{\text{core}}$ \\
Aristotle main table (T5) & refuted & 0 & (counterexample, no proof) \\
Aristotle main table (T4) & attempted & 0 & library-gap, no proof \\
Aristotle companion archive & closed & 3 & $\mathcal{A}_{\text{core}}$ \\
Aristotle companion archive & refuted & 1 & (counterexample, no proof) \\
Drafter closures (all) & closed & $\geq 1$ per closed row & $\mathcal{A}_{\text{core}}$ \\
\midrule
Release axiom-set union & n/a & n/a & $\mathcal{A}_{\text{core}}$ \\
\bottomrule
\end{tabular}}

\end{table*}

\paragraph{Why this matters.}
A capability benchmark that accepts proofs widening the axiom
set would confound tactic-search capability with axiom-cheat
pathology. Requiring axiom-set stability across all closed proofs
(enforced by the drafter sweep's closure criterion at
\S\ref{sec:methods}) makes the benchmark's closure counts
measurements of \emph{constructive} proof-search capability under
a fixed logical foundation.

\subsection{Failure-mode taxonomy of the drafter sweep}
\label{app:failure-modes}

Each drafter attempt resolves into one of five outcomes,
summarised in Table~\ref{tab:failure-modes}. \emph{Closed} attempts
type-check, satisfy the goal modulo no extra sorries, and do not
widen the axiom set beyond $\mathcal{A}_{\text{core}}$.
\emph{Drafter-wrote-sorry} attempts emit at least one
\texttt{sorry} token in the candidate proof body, signalling
explicit drafter abstention. \emph{Type-check failure} attempts
emit a candidate that the Lean elaborator rejects (unknown
identifier, type mismatch, unsolved goal, or tactic failure).
This is the dominant failure mode at $\sim 80\%$ of attempts and
is the conventional capability-blocked regime. \emph{Sample-side
error} attempts return no usable proof body (API timeout,
empty completion, format violation). These are sampling-pipeline
failures rather than capability signals. \emph{Alt-proof
non-close} attempts return a candidate distinct from the
reference Aristotle proof that nevertheless does not type-check.
This category is rare ($<0.5\%$).

\begin{table*}[!t]
\centering
\small
\caption{Per-drafter and aggregate failure-mode breakdown of
the $N=5$ sweep across T0/T1/T2 (T3 ablation columns in
Table~\ref{tab:ablation}). The dominant non-close mode is
type-check failure. Opus's sorry rate is $52/480 = 10.8\%$,
an order of magnitude above Sonnet's $2.4\%$ and the other
three drafters' $\leq 0.2\%$. Kimi and Mistral each have $85$
sample-side errors concentrated on T2 (truncated or empty API
completions), an integration artefact.}
\label{tab:failure-modes}
\begin{tabular}{lcccccc}
\toprule
Drafter & closed & sorry & type-fail & sample-err & alt-noclose & total \\
\midrule
Sonnet~4.6  & 50  & 14 & 530 & 0   & 1 & 595 \\
Kimi~K2.5   & 100 & 0  & 375 & 85  & 5 & 565 \\
Gemini~3 Pro & 103 & 1  & 461 & 0   & 0 & 565 \\
Mistral L~3 & 8   & 1  & 471 & 85  & 0 & 565 \\
Opus~4.6    & 61  & 52 & 365 & 0   & 2 & 480 \\
\midrule
aggregate   & 322 (11.6\%) & 68 (2.5\%) & 2202 (79.5\%) & 170 (6.1\%) & 8 (0.3\%) & 2770 \\
\bottomrule
\end{tabular}

\end{table*}

\subsection{Latency and cost profile}
\label{app:cost}

Table~\ref{tab:cost} reports per-drafter wall-clock and
per-attempt API cost aggregated across the $48$-slate
drafter sweeps (one greedy attempt at $T=0$ plus four
$T=0.7$ independent samples, per target). The cost column
is the per-attempt USD cost at Anthropic/Gemini/Mistral list
prices. Wall-clock is the full round-trip including network
latency.

\begin{table*}[!t]
\centering
\small
\caption{Drafter sweep latency and cost profile on the
$48$-slate. Attempt counts include duplicate runs after a
tokenizer fix. Per-attempt cost uses $T=0$ greedy plus $T=0.7$
samples at list prices. Total API-priced cost across the five
closed-source drafters is $\$37.30$. DSPv2 and Goedel-V2 run
locally with no API cost. Token counts are approximate.}
\label{tab:cost}
\resizebox{\linewidth}{!}{%
\begin{tabular}{lrrrrrr}
\toprule
Drafter & attempts & per-att wall & per-att cost & tok-in & tok-out & total \\
        &          & (sec)        & (USD)        & (avg)  & (avg)   & (USD) \\
\midrule
Claude Sonnet~4.6 & 821 & 7.9 & \$0.0070 & 160 & 432 & \$5.71 \\
Kimi K2.5 & 480 & 23.1 & \$0.0048 & 157 & 4400 & \$2.30 \\
Gemini~3 Pro & 480 & 29.3 & \$0.0306 & 317 & 3017 & \$14.67 \\
Mistral Large~3 & 506 & 3.2 & \$0.0017 & 166 & 233 & \$0.87 \\
Claude Opus~4.6 & 480 & 8.2 & \$0.0287 & 203 & 341 & \$13.75 \\
DSPv2-$7$B (community GPTQ-int8) & 240 & 3.2 & \$0 & $\sim$190 & $\sim$130 & \$0 \\
Goedel-V2-$8$B (community GGUF Q$6$\_K) & 240 & 21.8 & \$0 & 111 & 535 & \$0 \\
\bottomrule
\end{tabular}}

\end{table*}

\paragraph{Cost per closed target.}
We report API cost per closed target, defined as the
per-drafter total API cost on the $48$-slate divided by the
number of targets closed at pass@$5$. This avoids treating
attempts as independent Bernoulli trials, which they are not.
On the $48$-slate at pass@$5$: Mistral closes $2/48$ at
$\$0.87$ total ($\$0.44$ per closed target on the neutral
prompt). Sonnet closes $18/48$ at $\$5.71$ total ($\$0.32$
per closed target). Kimi closes $18/48$ at $\$2.30$ total
($\$0.13$ per closed target). Gemini closes $22/48$ at
$\$14.67$ total ($\$0.67$ per closed target). Opus closes
$13/48$ at $\$13.75$ total ($\$1.06$ per closed target).
DSPv2 closes $12/48$ at $\$0$ marginal API cost. The ranking changes if measured
per attempt rather than per closed target, because attempt
cost and closure are not linearly related. We do not derive a
per-attempt expected closure rate.

\subsection{Glossary of formal-methods terms}
\label{app:glossary}

\begin{description}[leftmargin=1.3em, itemsep=0.2em]
\item[Anytime-valid CS.] A confidence sequence
$(C_t)_{t\geq 1}$ at level $1-\alpha$ is \emph{anytime-valid} when
$\Prob(\forall t \geq 1 : \mu \in C_t) \geq 1-\alpha$. The
guarantee holds at every stopping time simultaneously, fixing
optional stopping in sequential experiments.
\item[Rate function $\eta_F$.] Per-family CS half-width
function. Smaller $\eta_F$ means a tighter CS. Cross-family
inequalities of the form $\eta_F \leq \eta_G$ rank the families.
\item[$\Phi$-substitution.] Log-domain substitution
$x \mapsto \exp(\lambda^\star x - (\lambda^\star)^2 \sigma^2 / 2)$
that converts an additive cumulant inequality to a multiplicative
exponential supermartingale. Central to the cross-family ranking
proofs.
\item[Lateral re-encoding.] Capability of proposing the
$\Phi$-substitution autonomously from the goal statement without
explicit prompt guidance.
\item[Tactic vocabulary.] Capability of selecting the
sequence of Mathlib lemmas (\texttt{Real.sqrt\_le\_sqrt},
\texttt{nlinarith}, etc.) that closes a re-encoded goal.
\item[Axiom-clean closure.] A Lean kernel certificate where
\texttt{\#print axioms} returns exactly
$\{\texttt{propext}, \texttt{Classical.choice}, \texttt{Quot.sound}\}$.
The canonical Mathlib core.
\item[MCTS.] Monte Carlo Tree Search, a search policy used by
some specialised theorem provers to explore tactic sequences.
Aristotle's internal search policy is not publicly documented.
\item[Optional stopping.] Choosing the sample size based on
the data already observed. Inflates classical fixed-horizon
type-I error from the nominal $\alpha$ to several times higher
in practice.
\item[pass@$N$.] A target is closed at pass@$N$ if at least one
of $N$ independent attempts produces an axiom-clean Lean kernel
certificate.
\end{description}

\subsection{How to extend the benchmark}
\label{app:how-to-extend}

The benchmark's tier structure is designed to admit new targets
at any of T0 through T3, plus the T4/T5 frontier. To submit a
new target:

\begin{enumerate}[leftmargin=1.3em, itemsep=0.15em]
\item Author the theorem statement in a Lean module under
\texttt{formal-avs-lean/FormalAvs/}, with the appropriate
imports.
\item Assign a paper tier: T0 (Mathlib-canon), T1
(single-function friendly), T2 (single-function challenging), T3
(cross-function), T4 (Mathlib-gap), T5 (capability-ceiling).
\item Add a JSONL row to the released \texttt{benchmark.jsonl}
with: \texttt{benchmark\_id}, \texttt{statement\_lean},
\texttt{tier}, \texttt{required\_imports},
\texttt{requires\_formal\_avs\_lean}, and the appropriate
\texttt{drafter\_close\_rate\_*} fields (initially \texttt{null}).
\item Run the drafter sweep at $N=5$ via the released
\texttt{sweep.py} against the four headline drafters and DSPv2.
Populate the \texttt{drafter\_close\_rate\_*} fields.
\item For T3 targets, additionally provide the \textsc{Sub} and
\textsc{Sub+Chain} hint columns and run the 3-column ablation.
\item For Aristotle coverage, submit the target as a session
via the public Aristotle API. Record session\_id, verdict, and
axiom audit in \texttt{aristotle\_history.jsonl}.
\item Open a pull request against the public benchmark
repository.
\end{enumerate}

The benchmark is versioned. Snapshot tags
(\texttt{formal-avs-v1.0}, \texttt{v1.1}, \ldots) freeze the
target set, drafter close-rate columns, and Aristotle history at a
specific commit so reported results stay reproducible.

\subsection{Aristotle session metadata}
\label{app:aristotle-session-metadata}

The 14-session companion archive records, per session,
fingerprint, target-set, verdict, axiom-clean status, and
elapsed wall time. Table~\ref{tab:aristotle-sessions} summarises.

\begin{table*}[!t]
\centering
\small
\caption{Aristotle 14-session archive summary. All
``closed'' rows are axiom-clean
($\mathcal{A}_{\text{core}}$).
\textsuperscript{$\dagger$}~companion-archive sessions
(\S\ref{app:aristotle-companion}).
\textsuperscript{$\ddagger$}~library-development sessions whose
targets are not in the benchmark. Full per-session JSON in
\texttt{aristotle\_history.jsonl}.}
\label{tab:aristotle-sessions}
\begin{tabular}{lp{3.5cm}p{4cm}r}
\toprule
Session id & Scope & Verdict & Wall (min) \\
\midrule
\texttt{74fe5c49} & 6 T3 cross-function & closed & $42$ \\
\texttt{f0bac40c} & 48-slate (T0/T1/T2) & 46 closed + 2 refuted & $51$ \\
\texttt{76c5d525} & 7 paper additions & closed & $36$ \\
\texttt{bd3ea302} & 4 T4 library-gap & no proof & $50$ \\
\texttt{5ac0658d} & 1 T5 capability-ceiling & refuted (symmetry) & $44$ \\
\texttt{d2755ea2} & Gaussian small-ball & refuted (symmetry) & $20$ \\
\texttt{54614669} & corrected small-ball & closed & $14$ \\
companion\textsuperscript{$\dagger$} (3 sessions) & sharp-rank, entropy, Kelly & all closed & $\sim 30$ each \\
4 additional\textsuperscript{$\ddagger$} & library-development & mixed & varies \\
\bottomrule
\end{tabular}

\end{table*}

\subsection{T0 stripped-prompt disambiguation}
\label{app:t0-stripped-prompt}

The headline T0 close-rate of $0/28$ across the four generalist
drafters could plausibly be explained either by a tactic-prior or
retrieval gap on Mathlib-canon lemmas, or by an artefact of the
hint-list prompt template (the banned-tactics clause, the
allowed-tactics enumeration, or the namespace preamble suppressing
otherwise-feasible closures). To disambiguate, we pre-registered and
ran a stripped-prompt control: theorem statement only, no
banned-tactics list, no allowed-tactics enumeration, no namespace
preamble, on $6$ of the $7$ T0 targets across the four drafters at
$N=5$ ($120$ attempts total). The stripped prompt closed $0/120$,
matching the hint-list rate. This rules out the prompt-wrapping
hypothesis. The remaining T0 target
(\texttt{div\_le\_div\_of\_nonneg\_left}) was excluded because
its as-stated form was refuted by Aristotle
(\S\ref{sec:aristotle-refute}). The corrected form is a fair test
for a future revision.

\subsection{N=25 T3 replicate}
\label{app:n25-t3-replicate}

To shrink the wide 95\% Wilson CI on the headline T3 aggregate
$2/24$, we ran a replicate sweep at $N=25$ on the six T3
cross-function targets across the four generalist drafters
(Sonnet~4.6, Kimi~K2.5, Gemini~3 Pro, Mistral Large~3) for a
total of $4 \times 6 \times 25 = 600$ attempts. Five of the six
T3 targets reconstruct cleanly into the sweep driver. One
(\texttt{etahr\_over\_etabetting\_gt\_one}) fails statement
reconstruction at $N=25$ and is excluded.

The unit of analysis for pass@$N$ is the target-drafter cell
(\S\ref{sec:methods}, drafter sweep protocol), not the
individual attempt. We therefore report cell-level pass@$25$
rather than aggregating attempts as if independent.

\begin{table*}[!t]
\centering
\small
\begin{tabular}{lc}
\toprule
Drafter & T3 pass@$25$ (5 targets, neutral prompt) \\
\midrule
Claude Sonnet 4.6 & $0/5$ \\
Kimi K2.5 & $0/5$ \\
Gemini 3 Pro & $0/5$ \\
Mistral Large 3 & $0/5$ \\
\midrule
Aggregate (4 drafters $\times$ 5 targets) & $0/20$ cells \\
\bottomrule
\end{tabular}
\end{table*}

\smallskip
The cell-level Wilson 95\% upper bound on pass@$25$ at $0/20$
is approximately $16.1\%$. This is evidence on the no-hint
regime on $5$ of the $6$ T3 targets. It does not directly bound
the headline hint-list result ($2/24$ at $N=5$ across $6$ T3
targets), since the prompts and target sets differ. The total
compute spent on this replicate ($4 \times 5 \times 25 = 500$
attempts) is reported under \S\ref{app:cost} for transparency,
but is not used to compute an attempt-level pass rate.

\section{Prompt-sensitivity detail}
\label{app:prompt-sensitivity-detail}

The hint-list prompt improves three of four generalist drafters and
reduces Mistral's closure rate.

\textbf{Mistral Large~3.} The hint-list acts as an action-space
restriction. Mistral's hint-list attempts converge to a deterministic
\texttt{:= by positivity} pattern, selecting one tactic from the
enumerated list ($0/48$ on hint-list). Without the enumeration,
pass@$5$ rises to $2/48$, recovering the T3 $\sqrt{2}$ identity
and one T1 target.

\textbf{Sonnet, Kimi, Gemini.} The hint-list provides a useful
prior. Removing it reduces pass@$5$ by a drafter-specific factor:
Sonnet $18 \to 3$ ($6\times$), Kimi $18 \to 13$ ($1.4\times$),
Gemini $22 \to 13$ ($1.7\times$). Sonnet's drop is not a parsing
artifact: $88.6\%$ of neutral attempts contain a parseable Lean
block. The prompt-sensitivity result argues
against reading the headline table as a single leaderboard.

\section{Inference-format analysis}
\label{app:inference-format}

On our slate, DSPv2-7B~\cite{dspv22025} closes $9/48$ at pass@$5$
while Goedel-Prover-V2-8B~\cite{goedelprover2025} closes $5/48$
under single-shot, despite Goedel-V2 reporting higher closure rates
on published benchmarks (MiniF2F, ProofNet). Four of five Goedel-V2
closures overlap with DSPv2. The gap is a hardware-constraint
artifact: DSPv2 outputs pure Lean in $50$--$80$ tokens (autocomplete
format), while Goedel-V2 is chat-trained with a recommended
$32$k-token inference window. On a $10$ GB consumer GPU the maximum
context is $4096$ tokens ($12.5\%$ of the recommended budget).

Under agentic mode (\texttt{cycle\_prove} $r{=}4$), Goedel-V2's
verifier-guided self-correction training activates: the model uses
Lean compiler feedback to iteratively revise its proofs, rising from
$5/48$ to $13/48$ ($2.6\times$). This surpasses DSPv2's single-shot
$9/48$, making Goedel-V2 the strongest local prover on our slate at
\$0 marginal cost. The self-correction capability is the
differentiator: DSPv2 (autocomplete-trained, no self-correction)
cannot benefit from a refinement loop.
